\documentclass[11pt]{article}
\usepackage[margin=1in]{geometry}
\usepackage{amsmath, amssymb, amsthm}
\usepackage{hyperref}

\newtheorem*{claim}{Claim}

\title{MAT 3150 --- Homework 2}
\author{}
\date{Due September 17}

\begin{document}
\maketitle

% ============================================================
\section*{Problem 1 (10 points)}
Prove that if $x$ and $y$ are arbitrary real numbers such that $x<y$,
there exists at least one irrational number $z$ satisfying $x<z<y$.

\begin{proof}
By the density of $\mathbb{Q}$ in $\mathbb{R}$, we can find a rational number $r_1$ such that $x < r_1 < y$. Applying density again to the interval $(r_1, y)$, we can find a rational number $r_2$ such that $r_1 < r_2 < y$.

Now consider
\[
z = r_1 + \frac{r_2 - r_1}{\sqrt{2}}.
\]

\textbf{$z$ is irrational.} Suppose, for the sake of contradiction, that $z$ is rational. Then
\[
z - r_1 = \frac{r_2 - r_1}{\sqrt{2}}
\]
has rational left-hand side (since $z$ and $r_1$ are both rational). So $\frac{r_2-r_1}{\sqrt2}$ is rational, say equal to $q$. Since $r_2 - r_1 \neq 0$ (as $r_1 < r_2$), we get $\sqrt{2} = \frac{r_2-r_1}{q}$, a ratio of two rationals, so $\sqrt{2}$ would be rational. This contradicts the irrationality of $\sqrt{2}$. Hence $z$ is irrational.

\textbf{$x < z < y$.} Since $\sqrt{2} > 1$ and $r_2 - r_1 > 0$,
\[
0 < \frac{r_2 - r_1}{\sqrt{2}} < r_2 - r_1.
\]
Adding $r_1$ throughout gives $r_1 < z < r_2$. Combined with $x < r_1$ and $r_2 < y$, this yields
\[
x < r_1 < z < r_2 < y,
\]
so in particular $x < z < y$.

\end{proof}

% ============================================================
\section*{Problem 2 (10 points)}
Find the supremum and the infimum of the sets
\[
  S_1 = \left\{ \frac{1}{n} : n \in \mathbb{N} \right\}
  \qquad \text{and} \qquad
  S_2 = \left\{ \frac{(-1)^n}{n} : n \in \mathbb{N} \right\}.
\]

\subsection*{$S_1$}
% Claim: sup S_1 = ___, inf S_1 = ___.
% For sup: identify the largest element explicitly (is it attained?).
% For inf: show 0 is a lower bound, then show it's the GREATEST lower
%   bound using the Archimedean property (for any eps > 0, find n with
%   1/n < eps).
By using the Archimedean property, we can show that $\inf S_1 = 0$. For any $\epsilon > 0$, there exists $n \in \mathbb{N}$ such that $1/n < \epsilon$, which means that $0$ is indeed the greatest lower bound of $S_1$.
The supremum of $S_1$ is $\sup S_1 = 1$, since the largest element of the set is $1$ (attained when $n=1$).

\subsection*{$S_2$}
For even $n$ the term $(-1)^n/n = 1/n$ is positive, and for odd $n$ the term is $-1/n$, negative. So the candidate supremum comes from the even terms and the candidate infimum from the odd terms.

\textbf{$\sup S_2 = \tfrac12$, attained at $n=2$.} Among the even terms $1/n$, the largest occurs at the smallest even $n$, namely $n=2$, giving $1/2$; every other even term satisfies $1/n \le 1/2$ for $n \ge 2$. Every odd term is negative, hence $< 1/2$. So $1/2$ is an upper bound for $S_2$, and since $1/2 \in S_2$ it is the maximum, hence $\sup S_2 = 1/2$.

\textbf{$\inf S_2 = -1$, attained at $n=1$.} Among the odd terms $-1/n$, the smallest (most negative) occurs at the smallest odd $n$, namely $n=1$, giving $-1$; every other odd term satisfies $-1/n \ge -1$ for $n\ge 1$. Every even term is positive, hence $> -1$. So $-1$ is a lower bound for $S_2$, and since $-1\in S_2$ it is the minimum, hence $\inf S_2 = -1$.

% ============================================================
\section*{Problem 3 (10 points)}
If $T,S$ are nonempty sets so that $s \le t$ for all $s\in S$ and $t\in T$,
prove that $\sup S$ and $\inf T$ exist, and prove the inequality
$\sup S \le \inf T$.

\begin{proof}
\textbf{Step 1: $\sup S$ exists.} $S$ is nonempty by hypothesis. Since $T$ is nonempty, fix some $t_0 \in T$. By hypothesis $s \le t_0$ for every $s \in S$, so $t_0$ is an upper bound for $S$. Thus $S$ is nonempty and bounded above, so by the completeness axiom, $\sup S$ exists.

\textbf{Step 2: $\inf T$ exists.} $T$ is nonempty by hypothesis. Since $S$ is nonempty, fix some $s_0 \in S$. By hypothesis $s_0 \le t$ for every $t \in T$, so $s_0$ is a lower bound for $T$. Thus $T$ is nonempty and bounded below, so by the completeness axiom, $\inf T$ exists.

\textbf{Step 3: $\sup S \le \inf T$.} Fix an arbitrary $t \in T$. Since $s \le t$ for every $s \in S$, $t$ is an upper bound for $S$, so by the least upper bound property $\sup S \le t$. This holds for every $t \in T$, so $\sup S$ is a lower bound for $T$. Since $\inf T$ is the greatest lower bound of $T$, it follows that $\sup S \le \inf T$.

\end{proof}

% ============================================================
\section*{Problem 4 (10 points)}
For $A$ and $B$ nonempty bounded subsets of real numbers, let $A+B$ be
the set of all sums $a+b$, where $a\in A$ and $b\in B$. Prove that
\[
  \sup(A+B) = \sup A + \sup B.
\]

\begin{proof}
Let $\alpha = \sup A$ and $\beta = \sup B$; these exist by the completeness axiom since $A$ and $B$ are nonempty and bounded above.

\textbf{Step 1: $\alpha+\beta$ is an upper bound for $A+B$.} Let $a\in A$ and $b\in B$ be arbitrary. Since $\alpha=\sup A$, $a\le\alpha$; since $\beta=\sup B$, $b\le\beta$. Adding, $a+b \le \alpha+\beta$. As $a,b$ were arbitrary, every element of $A+B$ is $\le \alpha+\beta$, so $\alpha+\beta$ is an upper bound for $A+B$. In particular, $A+B$ is nonempty and bounded above, so $\sup(A+B)$ exists and
\[
\sup(A+B) \le \alpha+\beta.
\]

\textbf{Step 2: $\alpha+\beta$ is the least upper bound.} Let $\varepsilon>0$ be arbitrary. Since $\alpha=\sup A$, $\alpha-\varepsilon/2$ is not an upper bound for $A$, so there exists $a\in A$ with $a > \alpha-\varepsilon/2$. Similarly, since $\beta=\sup B$, there exists $b\in B$ with $b>\beta-\varepsilon/2$. Then $a+b\in A+B$ and
\[
a+b > \left(\alpha-\frac{\varepsilon}{2}\right)+\left(\beta-\frac{\varepsilon}{2}\right) = \alpha+\beta-\varepsilon.
\]
Thus for every $\varepsilon>0$ there is an element of $A+B$ exceeding $\alpha+\beta-\varepsilon$, so no number strictly less than $\alpha+\beta$ can be an upper bound for $A+B$. Hence
\[
\sup(A+B) \ge \alpha+\beta.
\]

Combining Steps 1 and 2, $\sup(A+B) = \alpha+\beta = \sup A + \sup B$.

\end{proof}

\end{document}
